% ==============================================================================
% capitulos/02-estado-del-arte.tex - Estado del Arte
% ==============================================================================

\chapter{Estado del Arte}

\begin{flushleft}
Para el desarrollo de la presente plataforma, es fundamental realizar un análisis comparativo de las soluciones tecnológicas existentes en el mercado y en el ámbito institucional que abordan la gestión de quejas. Este estudio permite identificar las deficiencias funcionales de las herramientas actuales y definir el valor agregado de nuestra propuesta, especialmente en lo que respecta a la automatización y la inteligencia aplicada al derecho administrativo.
\end{flushleft}

\justifying

\section{Sistemas orientados a la denuncia y gestión de casos}
En la actualidad, existen diversas plataformas diseñadas para la captación de incidentes, quejas o denuncias, tanto en el sector educativo como en el ámbito laboral y civil. A continuación, se describen las herramientas más relevantes analizadas durante esta investigación:

\begin{itemize}
    \item \textbf{BullyButton.com:} Es una herramienta comercial orientada a instituciones educativas para el reporte de incidentes de acoso (\textit{bullying}). Su enfoque principal es la recolección rápida de datos y la generación de estadísticas para los administradores escolares \cite{7}.
    \item \textbf{Ventanilla Escolar (Gobierno de México):} Plataforma gubernamental diseñada para ofrecer orientación ciudadana sobre incidentes en planteles escolares. Funciona principalmente como un canal de comunicación y asesoría inicial \cite{8}.
    \item \textbf{Alerta IPN:} Herramienta de carácter institucional dentro del Instituto Politécnico Nacional utilizada para reportar sucesos anómalos o emergencias dentro de las instalaciones. Aunque es efectiva para alertas inmediatas, no cuenta con un flujo de seguimiento jurídico-administrativo \cite{9}.
    \item \textbf{UNAM - Denuncia Línea:} Sistema implementado por la Universidad Nacional Autónoma de México para el registro de quejas relacionadas con la normativa universitaria. Es uno de los referentes más cercanos en cuanto a estructura institucional \cite{10}.
    \item \textbf{PROFEDET Digital:} Plataforma de la Procuraduría Federal de la Defensa del Trabajo que utiliza herramientas digitales para la defensa de los derechos laborales, permitiendo un seguimiento más formal de los casos \cite{11}.
\end{itemize}

\section{Análisis Comparativo de Funcionalidades}
Con el objetivo de determinar el grado de innovación del presente proyecto, se realizó una comparativa técnica centrada en las capacidades de gestión operativa y la integración de tecnologías avanzadas.

\begin{table}[htbp]
\centering
\small
\begin{tabular}{|p{4cm}|c|c|c|c|c|}
\hline
\textbf{Característica} & \textbf{BullyButton} & \textbf{Ventanilla} & \textbf{Alerta IPN} & \textbf{UNAM} & \textbf{PROFEDET} \\
& \textbf{[7]} & \textbf{Escolar [8]} & \textbf{[9]} & \textbf{Denuncia [10]} & \textbf{Digital [11]} \\ \hline
Plataforma Web & Sí & Sí & Sí & Sí & Sí \\ \hline
Registro de quejas formal & Sí & Sí & No & Sí & Sí \\ \hline
Seguimiento por folio & Sí & No & No & No & No \\ \hline
Generación de estadísticas & Sí & No & No & Sí & Sí \\ \hline
Clasificador de casos con IA & No & No & No & No & No \\ \hline
\textit{Arquitectura Microservicios} & No & No & No & No & No \\ \hline
\end{tabular}
\caption{Comparativa de funciones de gestión operativa y técnica.}
\label{tab:comparativa_estado_arte}
\end{table}

\section{Diferenciación y Ventaja Competitiva}
Tras el análisis detallado en la Tabla \ref{tab:comparativa_estado_arte}, se concluye que la mayoría de las herramientas institucionales actuales operan como repositorios de información o canales de comunicación unidireccionales que requieren de una intensa intervención manual para la clasificación y gestión de incidentes. 

La propuesta desarrollada en este Trabajo Terminal presenta tres diferenciadores clave:

\begin{enumerate}
    \item \textbf{Automatización Cognitiva:} A diferencia de los sistemas analizados, esta plataforma incorpora un clasificador basado en Procesamiento de Lenguaje Natural (PLN) que asiste en la determinación de competencia desde el primer contacto, reduciendo significativamente los tiempos de respuesta.
    \item \textbf{Arquitectura de Grado Empresarial:} El uso de microservicios con Java Spring Boot garantiza que el sistema no sea un prototipo aislado, sino una infraestructura escalable y resiliente capaz de integrarse con otros sistemas institucionales del IPN.
    \item \textbf{Trazabilidad Total:} El sistema cierra el ciclo de coordinación entre el quejoso y el cuerpo de abogados mediante una experiencia unificada de seguimiento, eliminando la fragmentación de datos característica de los métodos tradicionales (archivo físico y herramientas de oficina descentralizadas).
\end{enumerate}

Este enfoque no solo moderniza la operación interna de la Defensoría de los Derechos Politécnicos, sino que establece un precedente tecnológico en la administración de justicia administrativa dentro del Instituto Politécnico Nacional.